We assume that requests arrive \emph{online}. A request $J_i$ and its
properties are known only when it arrives to the system.  We use
$a_i$ to denote its arrival time and $\ell_i$ to denote its
processing time. In the weighted case $J_i$ has a non-negative
weight $w_i$; the unweighted case corresponds to $w_i=1$ for all
$i$. Consider an online scheduling algorithm $A$. Let $f^A_i$ denote
the completion time or finish time of $J_i$ under $A$.  The response
time of $J_i$ under $A$ is $f^A_i - a_i$, in other words the total
duration spent by $J_i$ in the system before its processing is
finished. Now we define the delay factor metric. Here it is
assumed that each request $J_i$ has a \emph{deadline} $d_i$ that is
known upon its arrival. We refer to the quantity $S_i = (d_i-a_i)$
as the \emph{slack} of request $J_i$. If $f^A_i$ is the finish time
of $J_i$ under some schedule $A$ then its delay factor is defined as
$\max \{ 1, {(f^A_i - a_i)}/{(d_i - a_i)}\}$. Thus, delay factor
measures the ratio of the response time and the slack, unless the
request finishes by its deadline in which case we set its delay
factor to $1$.  In this paper we are interested in algorithms that
minimize the maximum response time and the maximum delay factor.
Given an online request sequence $\sigma$ and an algorithm $A$, let
$\alpha^A(\sigma) = \max_{J_i \in \sigma} \gamma^A_i$ where
$\gamma^A_i$ is either the response time or the delay factor of
$J_i$ in $A$; in the weighted case $\alpha^A(\sigma) = \max_{J_i \in
\sigma} w_i \gamma^A_i$.  We measure the performance of an algorithm
$A$ via worst-case competitive analysis: $A$ is $r$-competitive if
for all request sequences $\sigma$ we have $\alpha^A(\sigma) \le r
\alpha^*(\sigma)$ where $\alpha^*(\sigma)$ is the value of an
optimum offline schedule for $\sigma$.
